\documentclass[11pt]{article}

\setlength{\parindent}{0pt}
\usepackage{xltxtra}
\usepackage{hyperref}
\hypersetup{hidelinks}
\usepackage{url}
\usepackage{xcolor}
\definecolor{CVBlue}{RGB}{23,110,191}
\usepackage{calc}
\usepackage{graphicx}
\usepackage{tikz}
\usetikzlibrary{calc}
\usepackage{fontspec}
\usepackage{xeCJK}
\usepackage{enumitem}
\CJKsetecglue{} %% 取消中文与数字之间的间隙

%% 主文档字体设置
\setmainfont[
    Path = fonts/Main/,
    Extension = .otf,
    BoldFont = texgyretermes-bold.otf, % 加粗字体
]{texgyretermes-regular.otf} % 正文字体

% 中文字体设置
\setCJKmainfont[
    Path = fonts/hansans/,
    Extension = .ttf,
    BoldFont = NotoSansSC-Bold.ttf, % 加粗字体
]{NotoSansSC-Regular.otf} % 正文字体

\usepackage{relsize}
\usepackage{xspace}
\usepackage{fontawesome} 

% A4纸，上下左右边距
\usepackage[
    a4paper,
    left=1.2cm,
    right=1.2cm,
    top=1.5cm,
    bottom=1cm,
    nohead
]{geometry}

\renewcommand{\baselinestretch}{1.5} % 行间距设为1.5

\usepackage{titlesec}
\setlist{noitemsep} % 取消列表项间的额外间距
\setlist[itemize]{topsep=0.25em, leftmargin=*}
\setlist[enumerate]{topsep=0.25em, leftmargin=*}

% --- 用于控制【不同项目之间】的垂直距离 ---
\newlength{\interProjectSpacing}
\setlength{\interProjectSpacing}{0.9em} 
\newcommand{\projectsep}{\vspace{\interProjectSpacing}}

% --- 用于控制【项目标题】与下方【项目描述】的距离 ---
\newlength{\intraProjectTitleSep}
\setlength{\intraProjectTitleSep}{0.4em} 
\newcommand{\titlebreak}{\\[\intraProjectTitleSep]}

% --- 用于控制【项目描述】与下方【要点列表】的距离 ---
\newlength{\intraProjectListTopSep}
\setlength{\intraProjectListTopSep}{0.2em} 

% =======================================================================

\titleformat{\section}         
{\large\bfseries\raggedright} 
{}{0em}                      
{}                           
[{\color{CVBlue}\titlerule}]  
\titlespacing*{\section}{0cm}{*1.6}{*1.2}

\begin{document}
\pagenumbering{gobble}

%%%% 若本地有图片文件，可取消下方注释以显示照片和校徽
% \begin{tikzpicture}[remember picture, overlay] 
%     \node[anchor = north east] at ($(current page.north east)+(-2cm,-1.2cm)$) {\includegraphics[height=3cm]{avatar.jpg}};
% \end{tikzpicture}%
% \begin{tikzpicture}[remember picture, overlay] 
%     \node[anchor = north west] at ($(current page.north west)+(0.5cm,+1.0cm)$) {\includegraphics[height=6cm]{zju.png}};
% \end{tikzpicture}%

\centerline{\LARGE\bfseries{许珏凝}} 

\vspace{0.5em}
\centerline{\normalsize{\faPhone\ 15219616176 \quad \faEnvelopeO\ \href{mailto:3230105045@zju.edu.cn}{3230105045@zju.edu.cn}}} 

\vspace{1.5em}
    
\section{\makebox[\widthof{\faGraduationCap}][c]{\color{CVBlue}\faGraduationCap}\ 教育背景}    
\textbf{浙江大学} \hfill 2023.09 -- 至今\\[0.5em] 
数学科学学院 \quad 数学与应用数学专业 \quad 大三 \\
\textbf{学业成绩：} 前五学期总绩点 4.59/5.00 \quad \textbf{专业排名：} 4/69 \\
\textbf{外语水平：} CET-6 (635分)

\section{\makebox[\widthof{\faUsers}][c]{\color{CVBlue}\faUsers}\ 项目经历}

\textbf{美国大学生数学建模竞赛 (MCM/ICM) \quad S奖} \hfill 2025.01 -- 2025.02 \titlebreak
项目描述：针对奥运会奖牌榜预测与“名帅效应”量化问题，构建基于机器学习与统计推断的综合模型。
\begin{itemize}[nosep, topsep=\intraProjectListTopSep]
    \item \textbf{特征工程与奖牌预测}：基于运动员历史表现等10项特征构建 \textbf{随机森林 (Random Forest)} 模型，精准预测2028年奥运会奖牌榜 ($R^2=0.87$)；并结合 \textbf{K-means聚类与卡方检验}，深挖各参赛国的优势体育项目。
    \item \textbf{归因分析与多因素融合}：为剔除历史成绩干扰、孤立量化“名帅效应”，构建多元逻辑回归模型，并创新性引入 \textbf{优势分析 (Dominance Analysis)} 解决多变量依赖问题，证实教练因素对连续夺金的显著贡献度 (26.1\%)。
    \item \textbf{核心产出}：完成并提交25页全英文数据驱动数模论文，为国家奥委会的资源分配提供量化决策支持。
\end{itemize}

\projectsep

\textbf{代数曲线和模空间相关问题探讨 (院级SRTP项目) \quad 项目负责人} \hfill 2025.03 -- 至今 \titlebreak
项目描述：本项目旨在探索基础代数几何理论与高维数据降维、数值计算算法之间的交叉联系与应用表现。
\begin{itemize}[nosep, topsep=\intraProjectListTopSep]
    \item \textbf{理论基础与几何降维}：系统学习代数簇、层与上同调等核心理论；重点剖析 \textbf{Johnson-Lindenstrauss 引理}，探索高维数据随机投影后的几何结构保持性，研究其在数据科学降维算法中的底层数理逻辑。
    \item \textbf{数值模拟与计算验证}：跨越纯代数理论的抽象性，引入算法思维，规划并开展使用 \textbf{Python/SageMath} 对代数曲线族进行参数化数值模拟与动态绘制；通过椭圆曲线等显式计算实例，验证并可视化抽象模空间的拓扑性质。
\end{itemize}

\section{\makebox[\widthof{\faList}][c]{\color{CVBlue}\faList}\ 获奖情况}
\begin{itemize}[parsep=0.5ex]
    \item 2023-2024学年 国家奖学金
    \item 2024-2025学年 浙江省政府奖学金
    \item 浙江大学 一等奖学金
    \item 浙江大学 优秀学生
    \item 全国大学生数学竞赛（初赛） 二等奖
    \item 浙江省大学生数学竞赛（数学类） 三等奖
\end{itemize}
    
\section{\makebox[\widthof{\faCogs}][c]{\color{CVBlue}\faCogs}\ 技能与其他}
\begin{itemize}[nosep]
    \item \textbf{编程语言：} Python, C
    \item \textbf{排版与科学计算：} LaTeX, 了解数值模拟及算法求解工具
\end{itemize}
\end{document}
